\documentclass[12pt,a4paper]{ctexart}

% ===== 页面设置 =====
\usepackage[top=2.2cm,bottom=2.2cm,left=2.2cm,right=2.2cm,headheight=15pt]{geometry}
\usepackage{amsmath,amssymb}
\usepackage{enumitem}
\usepackage{xcolor}
\usepackage{tikz}
\usepackage{fancyhdr}

% ===== 页眉页脚 =====
\pagestyle{fancy}
\fancyhf{}
\fancyhead[L]{\textbf{高等数学（第七版·下册）\S11.2}}
\fancyhead[R]{\textbf{对坐标的曲线积分}}
\fancyfoot[C]{\thepage}
\renewcommand{\headrulewidth}{0.8pt}

\newcommand{\answerbox}[1]{%
  \begin{center}
  \fbox{\color{red}\large\bfseries #1}
  \end{center}
}

\begin{document}

\title{\Huge\bfseries 习题11-2~~三道路径对比计算}
\author{高等数学（同济大学数学系~编·第七版·下册）}
\date{\today}
\maketitle
\thispagestyle{fancy}

% ==================== 原题 ====================
\section*{题目}
\fbox{%
\parbox{\dimexpr\textwidth-2\fboxsep-2\fboxrule\relax}{%
计算对坐标的曲线积分：
\[
I = \int_L (xy-1)\,dx + x^2y\,dy
\]
其中 $L$ 分别为由点 $A(1,0)$ 到点 $B(0,2)$ 的下列弧段：
\begin{enumerate}[label=(\arabic*),leftmargin=*]
  \item 线段 $AB$：$2x + y = 2$
  \item 抛物线弧：$4x + y^2 = 4$
  \item 椭圆弧：$4x^2 + y^2 = 4$
\end{enumerate}
}%
}

\bigskip

\textbf{注：} 记 $P(x,y)=xy-1$，$Q(x,y)=x^2y$。检查是否与路径无关：
\[
\frac{\partial P}{\partial y} = x,\qquad
\frac{\partial Q}{\partial x} = 2xy,\qquad
\frac{\partial P}{\partial y} \neq \frac{\partial Q}{\partial x}.
\]
积分\textcolor{red}{与路径有关}，三条路径必须分别计算。

% ==================== (1) 线段 ====================
\section*{(1) 线段 $AB$：$2x + y = 2$}

\medskip

\noindent\textbf{参数化：}
直线 $2x+y=2$ 过 $A(1,0)$ 和 $B(0,2)$。
令 $x = t$，则 $y = 2-2t$。当 $t$ 从 $1 \to 0$ 时，点从 $A$ 移到 $B$。

\[
\begin{aligned}
x &= t,      &\quad dx &= dt,\\
y &= 2-2t,   &\quad dy &= -2\,dt,\\
t &: 1 \to 0.
\end{aligned}
\]

\noindent\textbf{代入积分：}
\[
\begin{aligned}
I_1 &= \int_1^0 \Bigl[t(2-2t)-1\Bigr]\,dt \;+\; t^2(2-2t)(-2)\,dt \\[4pt]
    &= \int_1^0 \Bigl[\,(2t-2t^2-1) \;-\; 2t^2(2-2t)\Bigr]\,dt \\[4pt]
    &= \int_1^0 \Bigl[\,2t-2t^2-1 \;-\; (4t^2-4t^3)\Bigr]\,dt \\[4pt]
    &= \int_1^0 \bigl(4t^3 - 6t^2 + 2t - 1\bigr)\,dt \\[4pt]
    &= \Bigl[\,t^4 - 2t^3 + t^2 - t\,\Bigr]_1^0 \\[4pt]
    &= (0) - (1-2+1-1) = 1.
\end{aligned}
\]

\answerbox{$I_1 = 1$}

% ==================== (2) 抛物线弧 ====================
\section*{(2) 抛物线弧：$4x + y^2 = 4$}

\medskip

\noindent\textbf{参数化：}
由 $4x + y^2 = 4$ 解出 $x = 1 - \dfrac{y^2}{4}$。
以 $y$ 为参数：令 $y = t$，则 $x = 1 - \dfrac{t^2}{4}$。

\[
\begin{aligned}
x &= 1 - \frac{t^2}{4}, &\quad dx &= -\frac{t}{2}\,dt,\\[4pt]
y &= t,                 &\quad dy &= dt,\\
t &: 0 \to 2 \quad (\text{从 }A\text{ 到 }B).
\end{aligned}
\]

\noindent\textbf{代入积分：}
\[
\begin{aligned}
I_2 &= \int_0^2 \Bigl[\Bigl(1-\frac{t^2}{4}\Bigr)t - 1\Bigr]\Bigl(-\frac{t}{2}\Bigr)\,dt
      \;+\; \Bigl(1-\frac{t^2}{4}\Bigr)^2 t \cdot dt \\[4pt]
    &= \int_0^2 \Bigl[\underbrace{\Bigl(t-\frac{t^3}{4}-1\Bigr)\Bigl(-\frac{t}{2}\Bigr)}_{\text{第1项}}
      \;+\; \underbrace{t\Bigl(1-\frac{t^2}{2}+\frac{t^4}{16}\Bigr)}_{\text{第2项}}\Bigr]\,dt \\[4pt]
    &= \int_0^2 \Bigl[\,\Bigl(-\frac{t^2}{2}+\frac{t^4}{8}+\frac{t}{2}\Bigr)
      \;+\; \Bigl(t-\frac{t^3}{2}+\frac{t^5}{16}\Bigr)\Bigr]\,dt \\[4pt]
    &= \int_0^2 \Bigl(\frac{t^5}{16} + \frac{t^4}{8} - \frac{t^3}{2} - \frac{t^2}{2} + \frac{3t}{2}\Bigr)\,dt \\[4pt]
    &= \Bigl[\,\frac{t^6}{96} + \frac{t^5}{40} - \frac{t^4}{8} - \frac{t^3}{6} + \frac{3t^2}{4}\,\Bigr]_0^2 \\[4pt]
    &= \frac{64}{96} + \frac{32}{40} - \frac{16}{8} - \frac{8}{6} + \frac{12}{4} \\[4pt]
    &= \frac{2}{3} + \frac{4}{5} - 2 - \frac{4}{3} + 3 \\[4pt]
    &= \Bigl(\frac{2}{3}-\frac{4}{3}\Bigr) + \frac{4}{5} + (3-2) \\[4pt]
    &= -\frac{2}{3} + \frac{4}{5} + 1 \\[4pt]
    &= \frac{-10 + 12 + 15}{15} = \frac{17}{15}.
\end{aligned}
\]

\answerbox{$I_2 = \dfrac{17}{15}$}

% ==================== (3) 椭圆弧 ====================
\section*{(3) 椭圆弧：$4x^2 + y^2 = 4$}

\medskip

\noindent\textbf{参数化：}
方程可写为 $\dfrac{x^2}{1} + \dfrac{y^2}{4} = 1$，是椭圆。用参数形式：
\[
x = \cos\theta,\qquad y = 2\sin\theta.
\]

验证：$4\cos^2\theta + 4\sin^2\theta = 4$ ✓

\[
\begin{aligned}
A(1,0) &: \cos\theta = 1,\; 2\sin\theta = 0 \;\Rightarrow\; \theta = 0,\\
B(0,2) &: \cos\theta = 0,\; 2\sin\theta = 2 \;\Rightarrow\; \theta = \frac{\pi}{2}.
\end{aligned}
\]
\[
dx = -\sin\theta\,d\theta,\qquad dy = 2\cos\theta\,d\theta,\qquad
\theta: 0 \to \frac{\pi}{2}.
\]

\noindent\textbf{代入积分：}
\[
\begin{aligned}
I_3 &= \int_0^{\frac{\pi}{2}}
      \Bigl[(\cos\theta \cdot 2\sin\theta - 1)(-\sin\theta)
      \;+\; \cos^2\theta \cdot 2\sin\theta \cdot 2\cos\theta\Bigr]\,d\theta \\[4pt]
    &= \int_0^{\frac{\pi}{2}}
      \Bigl[-(2\sin\theta\cos\theta - 1)\sin\theta
      \;+\; 4\cos^3\theta\sin\theta\Bigr]\,d\theta \\[4pt]
    &= \int_0^{\frac{\pi}{2}}
      \Bigl[-2\sin^2\theta\cos\theta + \sin\theta + 4\cos^3\theta\sin\theta\Bigr]\,d\theta \\[4pt]
    &= \underbrace{\int_0^{\frac{\pi}{2}}\sin\theta\,d\theta}_{\textcircled{1}}
       \;-\; 2\underbrace{\int_0^{\frac{\pi}{2}}\sin^2\theta\cos\theta\,d\theta}_{\textcircled{2}}
       \;+\; 4\underbrace{\int_0^{\frac{\pi}{2}}\cos^3\theta\sin\theta\,d\theta}_{\textcircled{3}}.
\end{aligned}
\]

分别计算三个积分：

\textcircled{1} $\displaystyle\int_0^{\frac{\pi}{2}}\sin\theta\,d\theta
                     = \bigl[-\cos\theta\bigr]_0^{\frac{\pi}{2}}
                     = 0 - (-1) = 1.$

\textcircled{2} 令 $u = \sin\theta$，$du = \cos\theta\,d\theta$：
   $\theta=0 \to u=0,\; \theta=\frac{\pi}{2} \to u=1$，
   $\displaystyle\int_0^1 u^2\,du = \frac{1}{3}.$

\textcircled{3} 令 $u = \cos\theta$，$du = -\sin\theta\,d\theta$：
   $\theta=0 \to u=1,\; \theta=\frac{\pi}{2} \to u=0$，
   $\displaystyle\int_1^0 u^3(-du) = \int_0^1 u^3\,du = \frac{1}{4}.$

\[
I_3 = 1 - 2 \cdot \frac{1}{3} + 4 \cdot \frac{1}{4}
     = 1 - \frac{2}{3} + 1
     = \frac{4}{3}.
\]

\answerbox{$I_3 = \dfrac{4}{3}$}

% ==================== 总结 ====================
\section*{总结与对比}

\begin{center}
\renewcommand{\arraystretch}{1.6}
\begin{tabular}{|c|c|c|}
  \hline
  \textbf{路径} & \textbf{方程} & \textbf{积分值 $I$} \\
  \hline
  (1) 线段      & $2x + y = 2$       & $1$ \\[4pt]
  \hline
  (2) 抛物线弧  & $4x + y^2 = 4$     & $\dfrac{17}{15} \approx 1.133$ \\[4pt]
  \hline
  (3) 椭圆弧    & $4x^2 + y^2 = 4$  & $\dfrac{4}{3} \approx 1.333$ \\[4pt]
  \hline
\end{tabular}
\end{center}

\bigskip

\textbf{核心认识：} 因为 $\dfrac{\partial P}{\partial y} \neq \dfrac{\partial Q}{\partial x}$，
该积分与路径有关，三条路径得出三个不同结果。
若被积函数满足 $\dfrac{\partial P}{\partial y} = \dfrac{\partial Q}{\partial x}$（如全微分 $d(x^2y)$），
则无论哪条路径结果都相同——这正是下一节格林公式要讨论的内容。

\end{document}
